

\begin{center}
    \subsection*{روابط و محاسبات}
\end{center}

\noindent
\textbf{
    داده های تجربی و اطلاعات به‌دست آمده از نمودار مربوط به آزمایش:
}

\begin{table}[htbp]
    \centering
    \renewcommand{\arraystretch}{1.2}
    \setlength{\tabcolsep}{14pt}
    \begin{latin}
        \begin{tabular}{ccccc}
            \toprule
                                                    & \rl{زرد} & \rl{سبز} & \rl{آبی} & \rl{بنفش} \\
            \midrule
            \rl{فرکانس نور} \((\qty{e14}{\hertz})\) & 5.19     & 5.49     & 6.88     & 7.41      \\
            \midrule
            \rl{ولتاژ معکوس} \((\unit{\volt})\)       & 0.42     & 0.55     & 1.04     & 1.30       \\
            \bottomrule
        \end{tabular}
    \end{latin}
    \caption{جدول داده های تجربی}
\end{table}


\vspace{-1.2cm}
\begin{flalign*}
    a \pm \Delta a & = 0.3847 \pm 0.0162 \ \unit{\volt\second} & \text{\RL{شیب نمودار}}  & \\
    b \pm \Delta b & = -1.574 \pm 0.102 \ \unit{\volt}         & \text{\RL{عرض از مبدأ}} &
\end{flalign*}

\vspace{-0.2cm}

\noindent
\textbf{نکته:}
خطای ولتاژ خوانده شده با مولتی‌متر \(\qty{0.01}{\volt}\) است.
خطای اصلی از تشخیص نقطه‌ای می‌آید که جریان صفر می‌شود و خطای کل ولتاژ حدود
\(\qty{0.03}{\volt}\)
برآورد شد.

\khat
\vspace{0.2cm}

\noindent
\textbf{
    روابط مربوط به تعیین ثابت پلانک، تابع کار فلز و فرکانس آستانه را بنویسید
    و هرکدام از پارامترها را در دستگاه \lr{SI} معرفی کنید
    و سپس مقدار عددی آنها را به‌دست آورید.
}


